\begin{questionS}
Implement a bounded partial stack using a monitor.  Give your implementation
the signature
\begin{scala}
class MonitorBoundedPartialStack[T](bound: Int) extends PartialStack[T]
\end{scala}
where |ParitalStack| is as in Exercise~\ref{ex:partialStack}.  The |push|
operation should block while the stack is full, i.e.~contains |bound|
elements.

Test your implementation using the linearizability testing framework.
\end{questionS}

%%%%%

\begin{answerS}
My code is below.  For convenience, I use a trait |PartialStack1| which gives
a default no-op implementation of |shutdown|.

The implementation is very similar to that for the bounded partial queue,
except adapted to a stack.  We use conditions |notFull| and |nonEmpty| to
signal that the relevant precondition holds.  Each operation waits on the
appropriate condition until the precondition holds.
%
\begin{scala}
trait PartialStack1[T] extends PartialStack[T]{
  def shutdown() = {}
}

/** A bounded partial stack implemented as a monitor. */
class MonitorBoundedPartialStack[T](bound: Int) extends PartialStack1[T]{
  /** The stack itself. */
  private val stack = new Stack[T]

  /** A monitor object, to control the synchronisations. */
  private val lock = new Lock

  /** Condition for signalling that the stack is not full. */
  private val notFull = lock.newCondition

  /** Condition for signalling that the stack is not empty. */
  private val notEmpty = lock.newCondition

  /** Push £x£.  Blocks while the queue is full. */
  def push(x: T) = lock.mutex{
    notFull.await(stack.length < bound) // Wait for a signal. (1)
    stack.push(x)
    notEmpty.signal() // Signal to a £pop£ at (2).
  }

  /** Dequeue a value.  Blocks until the queue is non-empty. */
  def pop(): T = lock.mutex{
    notEmpty.await(stack.nonEmpty) // Wait for a signal. (2)
    val result = stack.pop()
    notFull.signal() // Signal to a £push£ at (1).
    result
  }
} 
\end{scala}

The implementation can be tested precisely as in
Exercise~\ref{ex:partialStack}, except the push operation on the sequential
specification object needs to capture the precondition:
%
\begin{scala}
  def seqPush(x: Int)(st: SeqStack) : (Unit, SeqStack) = {
    require(st.length < bound); ((), st.push(x))
  }
\end{scala}
\end{answerS}
